% Generated draft: requires mathematical and source review.
\documentclass[11pt]{article}
\usepackage[margin=25mm]{geometry}
\usepackage{fontspec}
\setmainfont{Latin Modern Roman}
\usepackage{amsmath,amssymb,mathtools,unicode-math}
\setmathfont{Latin Modern Math}
\usepackage{enumitem,needspace,xurl,hyperref}
\setlength{\parindent}{0pt}
\setlength{\parskip}{5pt}
\setlength{\emergencystretch}{3em}
\newcommand{\N}{\mathbb{N}}
\newcommand{\Z}{\mathbb{Z}}
\newcommand{\Q}{\mathbb{Q}}
\newcommand{\R}{\mathbb{R}}
\newcommand{\C}{\mathbb{C}}
\newcommand{\F}{\mathbb{F}}
\newcommand{\A}{\mathbb{A}}
\newcommand{\PP}{\mathbb{P}}
\newcommand{\E}{\mathbb{E}}
\newcommand{\eps}{\varepsilon}
\newcommand{\dd}{\,\mathrm{d}}
\begin{document}
\section*{UnsolvedMath ID: 701215}
\textbf{Source problem code:} catalogue entry
\subsection{Definitions and mathematical statement}
Let \(n\in\N\) with \(n>2\), and let \(A\) be a ring with identity. For the usual associative-ring interpretation, write \(\mathrm{GL}_n(A)\) for the group of invertible matrices and define the elementary subgroup by
\[
E_n(A)=\bigl\langle I_n+a e_{ij}:a\in A,\ 1\leq i,j\leq n,\ i\neq j\bigr\rangle\leq\mathrm{GL}_n(A),
\]
where \(e_{ij}\) is the matrix unit. For each proper two-sided ideal \(J\subsetneq A\), let \(\rho_J:M_n(A)\to M_n(A/J)\) be entrywise reduction. Interpreting a scalar matrix as any matrix \(\lambda I_n\) with \(\lambda\in A/J\), the required nonscalarity condition is
\[
\forall J\subsetneq A\text{ a two-sided ideal},\quad
\forall\lambda\in A/J,\qquad \rho_J(g)\neq\lambda I_n.
\]
For a group \(H\), put
\[
\operatorname{Comm}(H)=\{xyx^{-1}y^{-1}:x,y\in H\}.
\]
The supplied question asks whether there exist \(n>2\), a ring \(A\) with identity, and \(g\in E_n(A)\) satisfying the displayed nonscalarity condition such that \(g\) is not a single commutator. The ambient group for the commutator is not specified. For an explicitly chosen ambient group \(H\) containing \(E_n(A)\), the question becomes
\[
\exists n\in\N,\ n>2,\ \exists A,\ \exists g\in E_n(A):\quad
\left(\forall J\subsetneq A\text{ a two-sided ideal},\ \forall\lambda\in A/J,\ \rho_J(g)\neq\lambda I_n\right)
\ \land\ g\notin\operatorname{Comm}(H).
\]
In particular, taking \(H=E_n(A)\) and taking \(H=\mathrm{GL}_n(A)\) give distinct formulations; the supplied record does not identify which is intended. Associativity, the two-sided interpretation of ideals, and the meaning of scalar matrices over a possibly noncommutative ring also require confirmation from the source conventions. No commutativity assumption on \(A\) is supplied.
\subsection{Short English statement}
L. Vaserstein asks whether some elementary matrix of size greater than two remains nonscalar in every proper quotient ring yet cannot be expressed as a single commutator. The intended ambient group for commutators and the ring conventions require clarification. The registry labels the problem open; this is a supplied status, not an independently verified assessment.
\subsection{Sources}
\begin{itemize}
\item Local UnsolvedMath registry, permanent ID \(701215\), exact source problem code \texttt{catalogue entry}, title \emph{Kourovka Notebook Problem 16.23}. The supplied record labels the problem open and attributes it to L. Vaserstein, proposed in 2006.
\item Supplied source attribution: \emph{The Kourovka Notebook (Unsolved Problems in Group Theory)}, 16th issue (2006), Problem 16.23; also identified in the record with the 21st issue edition (2026), \texttt{arXiv:1401.0300v48}. The collection description names E. I. Khukhro and V. D. Mazurov as editors. These bibliographic details are supplied attribution and have not been independently checked.
\end{itemize}
\end{document}
